\documentclass[11pt]{article}

\usepackage[margin=1in]{geometry}
\usepackage{amsmath,amssymb,amsthm}
\usepackage{graphicx}
\usepackage[hidelinks]{hyperref}
\usepackage{enumitem}

\setlist[itemize]{leftmargin=1.5em}
\setlength{\parindent}{0pt}
\setlength{\parskip}{0.7em}

\newtheorem*{theorem*}{Theorem}
\newtheorem*{corollary*}{Corollary}

\newcommand{\diam}{\operatorname{diam}}
\newcommand{\dist}{\operatorname{dist}}
\newcommand{\dimH}{\dim_{\mathrm H}}

\begin{document}

\begin{center}
{\Large Partial Result Packet}\\[0.25em]
{\large Endpoint-positive tree products for the distortion \(2\) Hausdorff-Dvoretzky question}\\[0.35em]
\textbf{Status:} partial result, likely valid
\end{center}

\section*{Source problem}

The source paper is:
\[
\text{arXiv:1106.0879, Manor Mendel and Assaf Naor,}
\]
\[
\text{\emph{Ultrametric subsets with large Hausdorff dimension}.}
\]
It appeared in \emph{Inventiones Mathematicae} 192 (2013), 1--54.\\
The local copy is included as \texttt{source\_paper.pdf}.

Question 1.3 of the source asks, for \(\alpha>0\) and \(D>1\), how large a
Hausdorff-dimensional subset must exist in every compact metric space of
Hausdorff dimension at least \(\alpha\) if the subset is required to embed
into Hilbert space with distortion \(D\). In the low-distortion regime the
source proves positive dimension at distortion \(2+\delta\), and proves that
no positive-dimensional theorem is possible for distortions strictly smaller
than \(2\).

\begin{center}
\includegraphics[width=0.98\linewidth]{figures/open_problem_crop.png}
\end{center}

The passage above, from page 6 of the source PDF, states that Question 1.3 at
distortion \(2\) remains open.

\section*{Result proved here}

This packet proves that the source's own \(G(n,1/2)\)-fractal lower-bound
examples are not counterexamples at the endpoint. In fact, those compact
spaces themselves embed into ultrametric spaces with distortion \(2\). Combined
with the source's lower bound for distortions strictly smaller than \(2\), this
shows that the random-graph fractals are exactly sharp at the threshold.

\section*{Idea of the Proof}

The source's lower-bound compacta are boundaries of trees whose branching
levels carry finite metrics. Two boundary points are separated at their first
different coordinate; their distance is the corresponding finite-level
distance, multiplied by the level scale. If we replace each finite-level metric
by the constant metric equal to its diameter, the resulting boundary distance
is an ultrametric. Therefore the only possible distortion loss is the aspect
ratio of the finite metrics used at the levels. The \(G(n,1/2)\) blocks have
only two nonzero distances, \(1\) and \(2\), hence aspect ratio \(2\).

\section*{Theorem}

\begin{theorem*}
Let \(\alpha>0\), let \(n_0=1\) and \(n_k>1\), and for each \(k\ge 1\) let
\(d_k\) be a metric on \(\{1,\ldots,n_k\}\) normalized so that
\[
\diam_{d_k}\{1,\ldots,n_k\}=1,
\qquad
\min_{i\ne j} d_k(i,j)\ge \eta
\]
for some fixed \(\eta\in(0,1]\). On
\[
Y=\prod_{k=1}^\infty \{1,\ldots,n_k\}
\]
define, as in Section 10.1 of the source,
\[
\rho_\alpha(x,y)
=
\frac{d_{k(x,y)}(x_{k(x,y)},y_{k(x,y)})}
{\prod_{i=0}^{k(x,y)-1} n_i^{1/\alpha}},
\]
where \(k(x,y)\) is the first coordinate at which \(x\) and \(y\) differ.
Then \((Y,\rho_\alpha)\) embeds into an ultrametric space with distortion at
most \(1/\eta\).
\end{theorem*}

\begin{proof}
Define
\[
u_\alpha(x,y)=
\begin{cases}
0, & x=y,\\[0.3em]
\left(\prod_{i=0}^{k(x,y)-1} n_i^{1/\alpha}\right)^{-1},& x\ne y.
\end{cases}
\]
This is the usual boundary ultrametric of the rooted tree:
for any three boundary points, the two largest among the three first-splitting
levels coincide, so \(u_\alpha(x,z)\le
\max\{u_\alpha(x,y),u_\alpha(y,z)\}\).

By the normalization of \(d_k\), whenever \(x\ne y\),
\[
\eta\,u_\alpha(x,y)
\le
\rho_\alpha(x,y)
\le
u_\alpha(x,y).
\]
Equivalently,
\[
\rho_\alpha(x,y)\le u_\alpha(x,y)\le \eta^{-1}\rho_\alpha(x,y).
\]
Thus the identity map from \((Y,\rho_\alpha)\) to \((Y,u_\alpha)\) is an
embedding into an ultrametric space with distortion at most \(1/\eta\).
\end{proof}

\begin{corollary*}
For the \(G(n,1/2)\)-fractal examples used in Section 10.3 of the source to
prove Theorem 1.8, the whole compact space embeds into an ultrametric space,
and hence into Hilbert space, with distortion at most \(2\).
\end{corollary*}

\begin{proof}
In Section 10.3, the source uses finite random-graph metrics \(W_n\) in which
the positive distances are \(1\) and \(2\). After normalizing the diameter to
one, the nonzero distances are \(1/2\) and \(1\), so \(\eta=1/2\) in the
theorem. Therefore the resulting tree-product compactum is \(2\)-equivalent
to an ultrametric.

The source's Lemma 10.2 and Theorem 10.5 compute the Hausdorff dimension of
these tree products, and the proof of Theorem 1.8 then shows that every
positive-dimensional subset of the \(G(n,1/2)\)-fractal incurs Hilbert
distortion at least \(2\). Since the whole space embeds into an ultrametric
with distortion at most \(2\), every subset embeds into Hilbert space with
distortion at most \(2\). Thus, for this family, the endpoint threshold is
attained exactly.
\end{proof}

\section*{Verification and limitations}

This is a partial result. It does not prove that every compact metric space of
positive Hausdorff dimension has a positive-dimensional subset embedding into
Hilbert space with distortion \(2\). It proves that one central attempted
counterexample route, namely the source's \(G(n,1/2)\)-fractal construction,
cannot disprove the endpoint statement.

The proof uses only the explicit tree-product metric from Section 10.1 of the
source and the two-valued nature of the \(G(n,1/2)\) metrics described in
Section 10.3. The bounded literature check found the later
Mendel--Naor paper \emph{Ultrametric skeletons} (arXiv:1112.3416), which
strengthens the source's measure skeleton theorem and derives Talagrand's
nonlinear Dvoretzky theorem, but did not find a later full solution of the
distortion \(2\) Hausdorff-dimensional endpoint.

\section*{References}

\begin{itemize}
\item M. Mendel and A. Naor, \emph{Ultrametric subsets with large Hausdorff
  dimension}, arXiv:1106.0879; Invent. Math. 192 (2013), 1--54.
\item M. Mendel and A. Naor, \emph{Ultrametric skeletons}, arXiv:1112.3416;
  Proc. Natl. Acad. Sci. USA 110 (2013), no. 48, 19256--19262.
\item Y. Bartal, N. Linial, M. Mendel, and A. Naor, \emph{On metric
  Ramsey-type phenomena}, Ann. of Math. 162 (2005), 643--709.
\end{itemize}

\end{document}
